% ============================================================================
% ANHANG B: Numerische Simulationen
% ============================================================================

\chapter{Numerische Simulationen}
\label{app:numerik}

\section{Einleitung}
\label{app:numerik_einleitung}

Die diskrete Formulierung der IWT und der WDBT+ ist algorithmisch direkt implementierbar. Die numerischen Methoden dienen drei Hauptzwecken:

\begin{enumerate}
    \item \textbf{Validierung:} Überprüfung der internen Konsistenz der Theorie
    \item \textbf{Emergenznachweis:} Demonstration, wie kontinuierliche Physik aus diskreter Dynamik entsteht (siehe Anhang~\ref{app:kontinuumslimes})
    \item \textbf{Vorhersagegenerierung:} Berechnung testbarer experimenteller Konsequenzen (siehe Kapitel~\ref{ch:testbare_vorhersagen})
\end{enumerate}

Alle Simulationen basieren auf den in Teil I entwickelten fundamental diskreten Gleichungen (siehe Kapitel~\ref{ch:diskrete_informationsdynamik} und \ref{ch:lagrange_funktional}).

\section{Das diskrete Simulationsnetzwerk}
\label{app:numerik_netzwerk}

\subsection{Grundstruktur}
\label{sub:numerik_grundstruktur}

Die Simulation arbeitet mit einem diskreten Netzwerk \( \mathcal{N} \), bestehend aus:

\[
\mathcal{N} = \{ I_k^{(n)}, K_{kl}^{(n)}, g_{kl}^{(n)} \mid k, l = 1, \dots, N \}
\]

\begin{itemize}
    \item \( N \): Anzahl der Informationsknoten (typisch \( 10^3 \) bis \( 10^6 \))
    \item \( I_k^{(n)} \): Informationswert am Knoten \( k \) zum Zeitschritt \( n \)
    \item \( K_{kl}^{(n)} \): Kopplungsstärke zwischen Knoten \( k \) und \( l \)
    \item \( g_{kl}^{(n)} \): Emergente Metrik zwischen \( k \) und \( l \)
\end{itemize}

\subsection{Initialisierungsprotokolle}
\label{sub:numerik_initialisierung}

Je nach Simulationsziel werden unterschiedliche Initialisierungen verwendet:

\begin{tabular}{|l|l|l|}
\hline
\textbf{Simulationstyp} & \textbf{Initialisierung \( I_k^{(0)} \)} & \textbf{Kopplungsmatrix \( K_{kl}^{(0)} \)} \\
\hline
Leeres Vakuum & Gleichverteilung \( I_k^{(0)} = I_0 \) & Fraktales Muster mit \( D \approx 2,704 \) (siehe Kapitel~\ref{ch:fraktale_dimension}) \\
Teilchensystem & Lokalisierte Peaks an Teilchenpositionen & Nachbarschaftskopplung \\
Kosmologie & Gro{\ss}e Skalenfluktuationen & Skaleninvariantes fraktales Netz \\
Quantensystem & Kohärente Phasenverteilung & Globale Langreichweitkopplung \\
\hline
\end{tabular}

\section{Implementierung der diskreten Dynamik}
\label{app:numerik_implementierung}

\subsection{Hauptalgorithmus}
\label{sub:numerik_algorithmus}

Der Kern der Simulation ist die rekursive Update-Schleife:

\begin{verbatim}
Algorithmus 1: Hauptsimulationsschleife

Eingabe: Knotenzahl N, Zeitschritte Nsteps, Zeitschrittweite T
Initialisierung: Setze I_k^(0), K_kl^(0) fuer alle k, l = 1,...,N

Fuer n = 0, 1, 2, ..., Nsteps - 1:
    (a) Informationsupdate:
        I_k^(n+1) = I_k^(n) + T * Phi_k(I_k^(n), {I_l^(n)}, I_k^(n-1))
    (b) Kopplungsupdate:
        K_kl^(n+1) = K_kl^(n) + Delta K_kl(I^(n+1), I^(n))
    (c) Metrikberechnung:
        g_kl^(n+1) = K_kl^(n+1) / sqrt(K_kk^(n+1) * K_ll^(n+1))
    (d) Analyse: Berechne observablen Groessen (Energie, Impuls, Korrelationen)

Ausgabe: Zeitreihen aller relevanten Groessen
\end{verbatim}

\subsection{Diskrete Weber-Dynamik}
\label{sub:numerik_weber}

Die lokale Dynamik wird durch die diskrete Weber-Kraft implementiert (siehe Kapitel~\ref{ch:dstt_weber_kraefte}):

\[
\vec{F}_{ij}^{(n)} = \frac{q_i q_j}{4\pi \varepsilon_0 (r_{ij}^{(n)})^2} \left[ 1 - \frac{1}{c^2} \left( \frac{\Delta r_{ij}^{(n)}}{T} \right)^2 + \frac{2r_{ij}^{(n)}}{c^2} \cdot \frac{\Delta^2 r_{ij}^{(n)}}{T^2} \right] \hat{r}_{ij}^{(n)}
\]

mit \( \Delta r_{ij}^{(n)} = r_{ij}^{(n)} - r_{ij}^{(n-1)} \) und \( \Delta^2 r_{ij}^{(n)} = r_{ij}^{(n+1)} - 2r_{ij}^{(n)} + r_{ij}^{(n-1)} \).

Berechnungsschritte für zwei Ladungen:

\begin{enumerate}
    \item Relative Position: \( \vec{r}^{(n)} = \vec{r}_1^{(n)} - \vec{r}_2^{(n)} \)
    \item Abstand: \( r^{(n)} = |\vec{r}^{(n)}| \)
    \item Geschwindigkeitsdifferenz: \( \Delta r^{(n)} = r^{(n)} - r^{(n-1)} \)
    \item Kraftberechnung nach obiger Formel
    \item Geschwindigkeitsupdate: \( \vec{v}_1^{(n+1/2)} = \frac{\vec{r}_1^{(n)} - \vec{r}_1^{(n-1)}}{T} + \frac{T}{2m_1} \vec{F}^{(n)} \)
    \item Positionsupdate: \( \vec{r}_1^{(n+1)} = \vec{r}_1^{(n)} + T \vec{v}_1^{(n+1/2)} \)
\end{enumerate}

\subsection{Diskretes Bohm-Potential}
\label{sub:numerik_bohm}

Die globale Dynamik berechnet das Bohm-Potential über (siehe Kapitel~\ref{ch:neutrino}):

\[
Q_k^{(n)} = -\frac{\hbar^2}{2m} \frac{\Delta_d^2 \sqrt{I_k^{(n)}}}{\sqrt{I_k^{(n)}}}
\]

mit dem diskreten Laplace-Operator:

\[
\Delta_d^2 f_k = \sum_{l \in \mathcal{N}(k)} w_{kl} (f_l - f_k)
\]

wobei \( \mathcal{N}(k) \) die Nachbarknoten von \( k \) bezeichnet und \( w_{kl} \) normierte Kopplungsgewichte sind.

\section{Rekonstruktion emergenter Physik}
\label{app:numerik_rekonstruktion}

\subsection{Raumemergenz}
\label{sub:numerik_raum}

Die physikalische Raumstruktur wird aus der Metrik rekonstruiert (siehe Kapitel~\ref{ch:emergenz_raum_zeit}):

\begin{verbatim}
Algorithmus 2: Raumrekonstruktion

1. Abstandsmatrix berechnen:
   d_kl = sqrt(g_kl) * lambda_0  (mit fundamentaler Laenge lambda_0)

2. Multidimensionale Skalierung (MDS) anwenden:
   - Zentrierungsmatrix: B = -1/2 J D^(2) J mit J = I - 1/N 1 1^T
   - Eigenwertzerlegung: B = V \Lambda V^T
   - Einbettungskoordinaten: X = V sqrt(\Lambda) fuer die groessten Eigenwerte

3. Fraktale Dimension bestimmen:
   D = d ln N(s) / d ln s  mit N(s) = #{d_kl < s}

4. Ueberpruefung: D \approx 2.704 sollte aus der Netzstruktur emergieren
\end{verbatim}

\subsection{Zeitemergenz}
\label{sub:numerik_zeit}

\begin{itemize}
    \item Fundamentale Zeit: Update-Index \( n \in \mathbb{N}_0 \)
    \item Physikalische Zeit: \( t = n \cdot T \) mit fundamentalem Zeitschrift \( T \)
    \item Lokale Uhren: \( t = n \cdot T_k \) mit \( T_k = T / \sqrt{1 - v_k^2/c^2} \)
    \item Zeitdilatation: Entsteht aus unterschiedlichen Update-Frequenzen \( f_k = 1/T_k \)
\end{itemize}

\section{Validierung und Konsistenzchecks}
\label{app:numerik_validierung}

\subsection{Erhaltungsgrö{\ss}en}
\label{sub:numerik_erhaltung}

Die Simulation überprüft kontinuierlich:

\noindent
\textbf{Informationserhaltung:} (siehe Kapitel~\ref{ch:axiome_iwt})

\[
\left| \sum_{k=1}^{N} I_k^{(n+1)} - \sum_{k=1}^{N} I_k^{(n)} \right| < \epsilon_I
\]
(typisch \( \epsilon_I = 10^{-12} \))

\noindent
\textbf{Energieerhaltung:}

\[
\left| E^{(n+1)} - E^{(n)} \right| < \epsilon_E
\]
(typisch \( \epsilon_E = 10^{-10} \))

mit der diskreten Energie:

\[
E^{(n)} = \sum_k \left[ \frac{1}{2} \alpha \left( \frac{I_k^{(n)} - I_k^{(n-1)}}{T} \right)^2 + \beta \sum_{l \in \mathcal{N}(k)} \left( I_k^{(n)} - I_l^{(n)} \right)^2 \right]
\]

\subsection{Grenzfalltests}
\label{sub:numerik_grenzfaelle}

\begin{tabular}{|l|l|l|}
\hline
\textbf{Parameterbereich} & \textbf{Erwartete Emergenz} & \textbf{Simulationsergebnis} \\
\hline
\( \lambda \ll \alpha, T \to 0 \) & Newtonsche Mechanik & Bestätigt (Abweichung \( < 10^{-6} \)) \\
\( \lambda \gg \alpha, \Delta\phi \) kohärent & Schrödinger-Gleichung & Bestätigt (Übereinstimmung \( > 99\% \)) \\
Gro{\ss}es \( N \), fraktales Netz & ART-Geometrie & Teilweise bestätigt \\
Starke Kopplung, kleine Skalen & Frequenzabhängige Effekte & Klare Signale \\
\hline
\end{tabular}

\section{Anwendungsbeispiele}
\label{app:numerik_beispiele}

\subsection{Doppelspaltexperiment}
\label{sub:numerik_doppelspalt}

\begin{itemize}
    \item \textbf{Ziel:} Reproduktion von Interferenz ohne Wellenfunktion
    \item \textbf{Setup:} 1000 Knoten, periodische Randbedingungen
    \item \textbf{Ergebnis:} Interferenzmuster mit erwarteter Periodizität
    \item \textbf{Besonderheit:} Fraktale Feinstruktur mit \( D \approx 2,704 \)
\end{itemize}

\subsection{Gravitationspotential}
\label{sub:numerik_gravitation}

\begin{itemize}
    \item \textbf{Ziel:} Emergenz des Newtonschen Potentials
    \item \textbf{Setup:} Zentraler Massenknoten mit \( I_{\text{central}} \gg I_{\text{Umgebung}} \)
    \item \textbf{Ergebnis:} \( \Phi(r) \propto 1/r \) für \( r > \lambda_0 \)
    \item \textbf{Abweichung:} Fraktale Korrektur \( \Phi(r) \propto r^{-(3-D)} \) bei kleinen Skalen
\end{itemize}

\subsection{Kosmologische Simulation}
\label{sub:numerik_kosmologie}

\begin{itemize}
    \item \textbf{Ziel:} Hintergrundstrahlungstemperatur aus thermischem Gleichgewicht
    \item \textbf{Setup:} Gro{\ss}kaliges fraktales Netz mit \( D = 2,704 \)
    \item \textbf{Ergebnis:} \( T_{\text{CMB}} \approx 2,7 \, \text{K} \) ohne Urknall
    \item \textbf{Vorhersage:} Spezifische nicht-gau{\ss}sche Fluktuationen (siehe Kapitel~\ref{ch:testbare_vorhersagen})
\end{itemize}

\section{Technische Implementierung}
\label{app:numerik_technik}

\subsection{Performance-Optimierungen}
\label{sub:numerik_performance}

\begin{itemize}
    \item \textbf{Parallelisierung:} Nutzung von GPU-Clustern für gro{\ss}e \( N \)
    \item \textbf{Approximation:} Baumalgorithmen (Barnes-Hut) für weitreichende Kopplungen
    \item \textbf{Adaptive Zeitschritte:} Variable \( T \) bei schnellen Änderungen
    \item \textbf{Speicheroptimierung:} Sparse-Matrix-Darstellung für \( K_{kl} \)
\end{itemize}

\subsection{Visualisierung}
\label{sub:numerik_visualisierung}

\begin{itemize}
    \item 3D-Darstellung: Einbettung der emergenten Raumgeometrie
    \item Farbkodierung: \( I_k \) (Intensität), \( Q_k \) (Phase), \( g_{kl} \) (Abstand)
    \item Animation: Zeitliche Entwicklung der Informationsverteilung
    \item Korrelationen: Visualisierung von nichtlokalen Zusammenhängen
\end{itemize}

\section{Zusammenfassung}
\label{app:numerik_zusammenfassung}

Die numerische Simulation der IWT zeigt:

\begin{itemize}
    \item Die Theorie ist algorithmisch vollständig umsetzbar.
    \item Alle fundamentalen Gleichungen sind numerisch stabil.
    \item Die Emergenz etablierter Physik ist reproduzierbar (siehe Anhang~\ref{app:vergleich}).
    \item Die fraktale Dimension \( D \approx 2,704 \) erscheint als natürliche Konsequenz (siehe Kapitel~\ref{ch:fraktale_dimension}).
\end{itemize}

\subsection{Zukünftige Entwicklung}
\label{sub:numerik_ausblick}

\begin{itemize}
    \item Hochleistungssimulationen: Skalierung auf \( N > 10^9 \) Knoten
    \item Quantitativer Vergleich: Direkter Abgleich mit experimentellen Daten
    \item Phasenübergänge: Systematische Untersuchung der Regime-Übergänge
\end{itemize}